\subsection{Preprocessing Pipeline}

Raw dermoscopic images vary in size, illumination, and, because ISIC-2019 is
a union of several archives, in the presence of near-duplicate photographs of
the same lesion. The pipeline cleans all of this once, before any training, in
three steps.

\textbf{Near-duplicate removal.} We compute a 64-bit perceptual hash (pHash) for
every image and group images whose hashes are within a Hamming distance of 4
using union-find, keeping a single representative per group. This removed
1{,}263 near-duplicates, leaving 24{,}068 unique images. Removing duplicates
before splitting is what prevents the same lesion from leaking between training
and test.

\textbf{Lesion-grouped stratified split.} We split the data 70/10/20 into
training, validation, and test sets. The split is stratified jointly on class
and source (HAM/BCN/MSK) so each split has a representative mix, and it is
grouped by lesion identifier so that every image of a given physical lesion lands
in the same split. The result is 16{,}872 training, 2{,}400 validation, and
4{,}796 test images.

\textbf{Colour constancy and memmap.} Each kept image is corrected with
Shades-of-Gray colour constancy~\cite{finlayson2004shadesofgray} (Minkowski norm
$p=6$) to reduce the strong illumination differences between dermatoscopes,
then short-side-resized and centre-cropped to $224\times224$. The entire cleaned
dataset is written once to a single memory-mapped \texttt{uint8} array together
with the labels and metadata. At training time this array is opened read-only and
stays resident in the host's page cache, so after the first epoch a batch costs
little more than a memory copy, the key optimisation that makes single-worker
data loading on Windows fast enough.

\subsection{Overview of the Proposed Models}

We propose two models that share a common design: convolutional features, a
transformer trunk for global reasoning, and a cross-attention head that folds in
clinical metadata. They sit at two different points on the
accuracy/efficiency curve. Figure~\ref{fig:arch} sketches both. The first,
which we refer to as the \emph{lightweight hybrid}, is built to stay under 6M
parameters. The second, \emph{DEKAN}, spends more capacity in exchange for higher
balanced accuracy.

\begin{figure*}[!t]
\centering
\resizebox{\textwidth}{!}{%
\begin{tikzpicture}[
  font=\footnotesize,
  box/.style={rectangle, rounded corners=3pt, draw=black!65, align=center,
              minimum height=9mm, inner sep=4pt},
  main/.style={box, text width=3.2cm},
  side/.style={box, text width=2.8cm},
  title/.style={font=\footnotesize\bfseries, align=center},
  io/.style={main, fill=black!8},
  cnn/.style={main, fill=blue!10},
  tx/.style={main, fill=orange!18},
  meta/.style={side, fill=green!18},
  head/.style={main, fill=red!14},
  arr/.style={-{Latex[length=2.4mm,width=1.7mm]}, line width=0.65pt, black!75},
  metaarr/.style={arr, dashed, green!45!black},
  panel/.style={draw=black!25, rounded corners=4pt, inner sep=8pt},
  node distance=8mm and 10mm
]

% Panel A: Lightweight Hybrid
\node[title] (atitle) at (0,0) {(a) Lightweight Hybrid};
\node[io, below=7mm of atitle] (a0) {Input image\\$3\times224\times224$};
\node[cnn, below=of a0] (a1) {MobileNetV2 stem\\stride-16\\$96\times14\times14$};
\node[tx, below=of a1] (a2) {Patch projection\\49 tokens, $d{=}192$\\CLS $+$ position};
\node[tx, below=of a2] (a3) {ViT trunk\\6 blocks, 4 heads};
\node[meta, right=14mm of a3] (am) {Clinical metadata\\age, sex, site\\missing tokens};
\node[tx, below=of a3] (a4) {CLS image token};
\node[tx, below=of a4] (a5) {Metadata\\cross-attention};
\node[head, below=of a5] (a6) {Linear classifier\\8 classes};
\draw[arr] (a0) -- (a1);
\draw[arr] (a1) -- (a2);
\draw[arr] (a2) -- (a3);
\draw[arr] (a3) -- (a4);
\draw[arr] (a4) -- (a5);
\draw[arr] (a5) -- (a6);
\draw[metaarr] (am.west) -| ([xshift=5mm]a5.east) -- (a5.east);
\node[panel, fit=(atitle)(a0)(a1)(a2)(a3)(a4)(a5)(a6)(am)] {};

% Panel B: DEKAN
\node[title] (btitle) at (8.1,0) {(b) DEKAN};
\node[io, below=7mm of btitle] (b0) {Input image\\$3\times224\times224$};
\node[cnn, below left=8mm and 8mm of b0] (b1a) {DenseNet-121\\stem, stride-16};
\node[cnn, below right=8mm and 8mm of b0] (b1b) {EfficientNet-B0\\stem, stride-16};
\node[tx, below=19mm of b0] (b2) {Attention fusion\\learned query bank};
\node[tx, below=of b2] (b3) {TinyViT trunk\\8 blocks, 8 heads\\CLS $+$ position};
\node[meta, right=14mm of b3] (bm) {Clinical metadata\\age, sex, site\\missing tokens};
\node[tx, below=of b3] (b4) {CLS image token};
\node[tx, below=of b4] (b5) {Metadata\\cross-attention};
\node[head, below=of b5] (b6) {KAN classifier\\8 classes};
\draw[arr] (b0) -- (b1a);
\draw[arr] (b0) -- (b1b);
\draw[arr] (b1a) -- (b2);
\draw[arr] (b1b) -- (b2);
\draw[arr] (b2) -- (b3);
\draw[arr] (b3) -- (b4);
\draw[arr] (b4) -- (b5);
\draw[arr] (b5) -- (b6);
\draw[metaarr] (bm.west) -| ([xshift=5mm]b5.east) -- (b5.east);
\node[panel, fit=(btitle)(b0)(b1a)(b1b)(b2)(b3)(b4)(b5)(b6)(bm)] {};
\end{tikzpicture}%
}
\caption{High-level architecture of the two proposed models. Both share the same
philosophy convolutional features, a transformer trunk, and a metadata
cross-attention head at two capacity points. (a) the lightweight hybrid uses
a single truncated MobileNetV2 stem and a 6-layer ViT trunk with a linear head;
(b) DEKAN fuses two pretrained backbones (DenseNet-121, EfficientNet-B0) through
a learned attention bank, runs an 8-layer TinyViT trunk, and ends in a KAN head.}
\label{fig:arch}
\end{figure*}

\subsection{Lightweight Hybrid CNN--Transformer}

\subsubsection{CNN stem}

The stem is a MobileNetV2~\cite{sandler2018mobilenetv2} (width 1.0), loaded with
ImageNet weights and truncated at the stride-16 stage so that a
$224\times224$ image produces a $96\times14\times14$ feature map. Using a
pretrained convolutional front end gives the model generic edge and texture
detectors for free, which matters because the transformer that follows is trained
from scratch on only about 17{,}000 images. In our budget accounting the stem
accounts for roughly 1.81M of the model's parameters.

\subsubsection{Transformer trunk}

A $2\times2$ stride-2 convolution projects the $96$-channel, $14\times14$ feature
map into a sequence of $7\times7 = 49$ tokens of dimension 192; this both embeds
the patches and halves the spatial resolution so that self-attention, which is
quadratic in the number of tokens, stays cheap. A learnable class (CLS) token is
prepended and a learnable positional embedding added, giving 50 tokens. These
pass through six pre-norm transformer blocks with four attention heads, an MLP
expansion ratio of 2, and a linearly increasing stochastic-depth
rate~\cite{huang2016stochasticdepth} up to 0.1, which is what keeps a small
from-scratch transformer from overfitting at this data scale. The CLS token after
the final layer summarises the image globally.

\subsubsection{Metadata cross-attention head}

Clinical metadata enters through a single cross-attention block in which the
image CLS token is the query and three metadata tokens (age, sex, and
anatomical site) are the keys and values. Age, a continuous value, is mapped
to a token by a linear projection; sex and site, which are categorical, use
learned embeddings. Missing values are handled explicitly: age has a dedicated
learned ``missing'' vector that replaces its projection when the value is absent,
and the categorical embeddings reserve an extra slot for ``missing.'' Each token
also carries a learned type embedding so the attention can tell the three fields
apart. The block lets the model nudge the image representation using context,
for instance, raising the prior on actinic keratosis for an older patient,
without ever being forced to act on a fabricated value. The updated CLS token
feeds a linear classifier producing eight logits.

The complete lightweight hybrid has 3.98M parameters and 0.631 GMACs, comfortably
inside the budget of 6M parameters and 1 GMAC.

\subsection{DEKAN: Dual-Backbone Flagship Model}

DEKAN keeps the same overall shape but invests capacity where it buys accuracy.

\textbf{Two complementary backbones.} Instead of one stem it runs two in
parallel DenseNet-121~\cite{huang2017densenet} and
EfficientNet-B0~\cite{tan2019efficientnet}, both ImageNet-pretrained and both
truncated to the stride-16 stage. The two networks encode the same lesion through
different inductive biases, and the model is given the chance to use whichever is
more informative per region.

\textbf{Learned attention fusion.} Each backbone's feature map is projected to
49 tokens of dimension 256 and tagged with a per-backbone embedding. A learnable
bank of 49 query tokens then cross-attends over the concatenation of both token
sets (98 tokens) and produces a single fused sequence of 49 tokens. This is a
learned, content-dependent fusion rather than a fixed concatenation or average:
the query bank decides, position by position, how much to take from each
backbone.

\textbf{Transformer trunk and metadata.} A CLS token and positional embedding are
prepended to the fused sequence, which then passes through an eight-layer
TinyViT-style transformer trunk (dimension 256, eight heads, MLP ratio 4),
followed by the same metadata cross-attention head described above.

\textbf{KAN classifier.} The final classifier is a single Kolmogorov--Arnold
layer~\cite{liu2024kan} mapping the 256-dimensional CLS representation to eight
logits, with a grid size of 5 and cubic ($\text{order}=3$) B-splines. Its spline
basis is computed in full precision for numerical stability even under mixed
precision. We deliberately keep a linear-head version of the otherwise-identical
model so that the KAN's contribution can be measured directly (Chapter~5).

DEKAN has 16.45M parameters and 6.633 GMACs.

\subsection{Class-Balanced Focal Loss}

To counter the imbalance in Table~\ref{tab:classdist}, every model is trained
with Class-Balanced Focal Loss, which combines the effective-number reweighting
of Cui et al.~\cite{cui2019classbalanced} with the focal modulation of Lin et
al.~\cite{lin2017focal}. The effective number of samples for a class with $n$
examples is $E(n) = (1-\beta^{n})/(1-\beta)$, and the per-class weight is
$\alpha = (1-\beta)/(1-\beta^{n})$, normalised so the weights sum to the number
of classes. The per-sample loss is

\begin{equation}
\mathcal{L} = -\,\alpha_y\,(1 - p_y)^{\gamma}\,\mathrm{CE}_{\text{smooth}}(\text{logits}, y),
\end{equation}

where $p_y$ is the predicted probability of the true class, the
$(1-p_y)^{\gamma}$ term down-weights easy examples so that learning focuses on
the hard, often rare, cases, and $\mathrm{CE}_{\text{smooth}}$ is a
label-smoothed cross-entropy. In our configuration $\beta = 0.999$,
$\gamma = 2.5$, and the label-smoothing factor is $0.1$. The choice of
$\beta = 0.999$ rather than a more aggressive value is deliberate: a larger
$\beta$ produced an extreme weight ratio between the common and rare classes and
caused the rare classes to collapse early in training, whereas $0.999$ gives
stable optimisation.

\subsection{Training Protocol}

All models, including the five baselines and both proposed models, are trained with
an identical recipe so that comparisons are fair; only batch size, gradient
accumulation, and the backbone learning-rate scale differ, as dictated by each
model's memory footprint. The optimiser is AdamW~\cite{loshchilov2019adamw} with
a base learning rate of $3\times10^{-4}$, $\beta=(0.9, 0.999)$, and gradients
clipped to a maximum norm of 5. Pretrained CNN stems are fine-tuned gently at a
fraction of the base learning rate (0.3$\times$ for the hybrid, 0.1$\times$ for
DEKAN) while from-scratch components train at the full rate. The schedule is a
five-epoch linear warmup followed by cosine annealing~\cite{loshchilov2017sgdr}
to a minimum learning rate of $10^{-6}$. Training uses automatic mixed precision
and channels-last memory layout throughout, a weight exponential moving average
(EMA, decay $0.9998$) whose weights are used for validation and final
evaluation, and a weighted random sampler that draws rare classes more often.
Data augmentation combines horizontal and vertical flips, random resized crops,
colour jitter, RandAugment~\cite{cubuk2020randaugment}, and
Mixup~\cite{zhang2018mixup} (strength $\alpha=0.4$, applied to half the batches),
with metrics always computed on the original, unmixed labels. Table~\ref{tab:traincfg}
lists the per-model settings that vary.

\begin{table}[!t]
\centering
\caption{Per-model training settings. All other hyperparameters (loss, optimiser
base LR, schedule, EMA, augmentation) are identical across models.}
\label{tab:traincfg}
\begin{tabular}{lrrrrr}
\toprule
Model & Batch & Grad accum & Eff.\ batch & Epochs & Backbone LR \\
\midrule
ResNet-18           & 256 & 1 & 256 & 100 & 0.3$\times$ \\
MobileNetV2         & 256 & 1 & 256 & 100 & 0.3$\times$ \\
EfficientNet-B0     & 128 & 2 & 256 & 100 & 0.3$\times$ \\
MobileViT-S         & 128 & 2 & 256 & 150 & 0.3$\times$ \\
EfficientFormer-L1  & 128 & 1 & 128 & 100 & 0.3$\times$ \\
Hybrid (full)       & 96  & 3 & 288 & 150 & 0.3$\times$ \\
DEKAN (full)        & 64  & 4 & 256 & 150 & 0.1$\times$ \\
\bottomrule
\end{tabular}
\end{table}

\subsection{Evaluation Protocol}

The primary metric is balanced multi-class accuracy (BMA), the unweighted mean of
the eight per-class recalls, which gives every class common or rare equal
say. We additionally report macro-F1, macro-AUC, overall accuracy, and per-class
recall, and we visualise errors with a confusion matrix normalised by true class.
At inference we apply eight-view test-time augmentation (TTA): the model is run
on the original image plus a fixed set of flips and rotations and the softmax
outputs are averaged. The best checkpoint is selected by validation BMA and all
headline numbers are reported on the held-out test split.
